%Daten zur Institution

%\input{Dozentdaten}

%\renewcommand{\fachbereich}{Fachbereich}

%\renewcommand{\dozent}{Prof. Dr. . }

%Klausurdaten

\renewcommand{\klausurgebiet}{ }

\renewcommand{\klausurtyp}{ }

\renewcommand{\klausurdatum}{ . 20}

\klausurvorspann {\fachbereich} {\klausurdatum} {\dozent} {\klausurgebiet} {\klausurtyp}

%Daten für folgende Punktetabelle

\renewcommand{\aeins}{  3  }

\renewcommand{\azwei}{  3   }

\renewcommand{\adrei}{  0  }

\renewcommand{\avier}{  4  }

\renewcommand{\afuenf}{  5  }

\renewcommand{\asechs}{  4  }

\renewcommand{\asieben}{  9  }

\renewcommand{\aacht}{  0  }

\renewcommand{\aneun}{  8  }

\renewcommand{\azehn}{  6  }

\renewcommand{\aelf}{  0  }

\renewcommand{\azwoelf}{  3  }

\renewcommand{\adreizehn}{  3   }

\renewcommand{\avierzehn}{  0  }

\renewcommand{\afuenfzehn}{  0  }

\renewcommand{\asechzehn}{  7  }

\renewcommand{\asiebzehn}{  55  }

\renewcommand{\aachtzehn}{    }

\renewcommand{\aneunzehn}{    }

\renewcommand{\azwanzig}{    }

\renewcommand{\aeinundzwanzig}{    }

\renewcommand{\azweiundzwanzig}{    }

\renewcommand{\adreiundzwanzig}{    }

\renewcommand{\avierundzwanzig}{    }

\renewcommand{\afuenfundzwanzig}{    }

\renewcommand{\asechsundzwanzig}{    }

\punktetabellesechzehn

\klausurnote

\newpage

\setcounter{section}{K}

__NOEDITSECTION__

\inputaufgabeklausurloesung
{3}

{

Definiere die folgenden
\zusatzklammer {kursiv gedruckten} {} {} Begriffe.
\aufzaehlungsechs{Die
\stichwort {Gaußkrümmung} {}
zu einer differenzierbaren Fläche

\mavergleichskette
{\vergleichskette
{ Y
}
{ \subseteq }{ \R^3
}
{  }{
}
{    }{
}
{   }{
}
}
{}{}{}
in einem Punkt

\mavergleichskette
{\vergleichskette
{ P
}
{ \in }{ Y
}
{  }{
}
{    }{
}
{   }{
}
}
{}{}{.}

}{Eine
\stichwort {topologische Mannigfaltigkeit} {}
der Dimension $n$.

}{Ein
\stichwort {zeitunabhängiges Vektorfeld} {}
auf einer differenzierbaren Mannigfaltigkeit $M$.

}{Eine
\stichwort {Differentialform vom Grad $k$ } {}
auf einer differenzierbaren Mannigfaltigkeit $M$.

}{Eine \stichwort {riemannsche Mannigfaltigkeit} {.}

}{Ein
\stichwort {torsionsfreier} {}
linearer Zusammenhang auf dem Tangentialbündel einer riemannschen Mannigfaltigkeit $M$.
}

}
{

\aufzaehlungsechs{Man nennt das Produkt
\mathl{\kappa_1 \cdot \kappa_2}{} der beiden
\definitionsverweis {Hauptkrümmungen}{}{}
in $P$ die
\definitionswort {Gaußkrümmung}{}
der Fläche $Y$ in $P$.
}{Ein
\definitionsverweis {topologischer}{}{}
\definitionsverweis {Hausdorffraum}{}{}
$M$ heißt eine topologische Mannigfaltigkeit der Dimension $n$, wenn es eine
\definitionsverweis {offene Überdeckung}{}{}

\mavergleichskette
{\vergleichskette
{ M
}
{ = }{  \bigcup_{i \in I} U_i
}
{  }{
}
{    }{
}
{   }{
}
}
{}{}{}
derart gibt, dass jedes $U_i$
\definitionsverweis {homöomorph}{}{}
zu einer
\definitionsverweis {offenen Teilmenge}{}{}
des $\R^n$ ist.
}{Eine Abbildung
\maabbdisp {F} {M} {TM
} {}
mit der Eigenschaft, dass

\mavergleichskette
{\vergleichskette
{ F(P)
}
{ \in }{  T_PM
}
{  }{
}
{    }{
}
{   }{
}
}
{}{}{}
für jeden Punkt

\mavergleichskette
{\vergleichskette
{ P
}
{ \in }{ M
}
{  }{
}
{    }{
}
{   }{
}
}
{}{}{}
ist, heißt
\zusatzklammer {zeitunabhängiges} {} {}
Vektorfeld.
}{Eine
$k$-Differentialform
ist ein
\definitionsverweis {Schnitt}{}{}
im $k$-fachen
\definitionsverweis {Dachprodukt}{}{}
des
\definitionsverweis {Kotangentialbündels}{}{.}
}{Eine
\definitionsverweis {differenzierbare Mannigfaltigkeit}{}{}
$M$ heißt riemannsche Mannigfaltigkeit, wenn auf jedem
\definitionsverweis {Tangentialraum}{}{}

\mathbed {T_PM} {}
 {P \in M} {}
 {} {} {} {,}
ein
\definitionsverweis {Skalarprodukt}{}{}

\mathl{\left\langle - , -  \right\rangle_P}{} erklärt ist derart, dass für jede Karte
\maabbdisp {\alpha} { U } { V
} {}
mit
\mathl{V\subseteq \R^n}{} die Funktionen
\zusatzklammer {für \mathlk{1 \leq i,j \leq n}{}} {} {}
\maabbeledisp {g_{ij}} { V } {\R
} { Q } {g_{ij}(Q) =  \left\langle  T(\alpha^{-1})(e_i)  ,  T(\alpha^{-1}) (e_j)   \right\rangle_{ \alpha^{-1}(Q) }
} {,}
$C^1$-\definitionsverweis {differenzierbar}{}{}
sind.
}{Der Zusammenhang heißt
torsionsfrei,
wenn

\mavergleichskettedisp
{\vergleichskette
{ \nabla_V W- \nabla_W V
}
{ =} { [V,W]
}
{ } {
}
{ } {
}
{ } {
}
}
{}{}{}
für beliebige Vektorfelder $V,W$ gilt.
}

 }

__NOEDITSECTION__

\inputaufgabeklausurloesung
{3}

{

Formuliere die folgenden Sätze.
\aufzaehlungdrei{Der
\stichwort {Existenzsatz für den Paralleltransport auf einer Hyperfläche} {.}}{Der Satz über die lokale Beschreibung von Differentialformen.}{Die
\stichwort {Quaderversion} {}
des Satzes von Stokes.}

}
{

\aufzaehlungdrei{\faktsituation {Es sei

\mavergleichskette
{\vergleichskette
{ Y
}
{ \subseteq }{ W
}
{ \subseteq }{\R^n
}
{    }{
}
{   }{
}
}
{}{}{,}
$W$
\definitionsverweis {offen}{}{,}
eine
\definitionsverweis {differenzierbare Hyperfläche}{}{}
und sei
\maabb {\gamma} { I } { Y
} {}
eine
\definitionsverweis {differenzierbare Kurve}{}{.}
Es sei

\mavergleichskette
{\vergleichskette
{ t_0
}
{ \in }{ I
}
{  }{
}
{    }{
}
{   }{
}
}
{}{}{,}

\mavergleichskette
{\vergleichskette
{ \gamma(t_0)
}
{ = }{ P
}
{  }{
}
{    }{
}
{   }{
}
}
{}{}{}
und sei

\mavergleichskette
{\vergleichskette
{ v
}
{ \in }{ T_PY
}
{  }{
}
{    }{
}
{   }{
}
}
{}{}{}
ein
\definitionsverweis {Tangentialvektor}{}{.}}
\faktfolgerung {Dann gibt es ein eindeutig bestimmtes
\definitionsverweis {tangentiales Vektorfeld}{}{}
$F$ längs $\gamma$, das
\definitionsverweis {parallel}{}{}
ist und

\mavergleichskettedisp
{\vergleichskette
{ F(t_0)
}
{ =} { v
}
{ } {
}
{ } {
}
{ } {
}
}
{}{}{}
erfüllt.}
\faktzusatz {}
\faktzusatz {}}{Es sei $M$ eine
\definitionsverweis {differenzierbare Mannigfaltigkeit}{}{}
und
\mathl{U \subseteq M}{} eine
\definitionsverweis {offene Teilmenge}{}{}
mit einer Karte
\maabbdisp {\alpha} { U } { V
} {}
und
\mathl{V \subseteq \R^n}{} offen. Es seien
\maabbdisp {x_j} { U } {\R
} {} die zugehörigen Koordinatenfunktionen,
\mathl{1 \leq j \leq n}{.} Dann lässt sich jede auf $U$ definierte $k$-\definitionsverweis {Differentialform}{}{}

\mathl{\omega \in
{ \mathcal E }^{ k } (  U   )}{} eindeutig schreiben als

\mavergleichskettedisp
{\vergleichskette
{  \omega
}
{ =} { \sum_{J,\,  { \# \left(   J  \right) } = k}   f_J dx_J
}
{ } {
}
{ } {
}
{ } {
}
}
{}{}{}
mit eindeutig bestimmten Funktionen
\maabbdisp {f_J} { U } {\R
} {.}}{\faktsituation {Es sei

\mavergleichskette
{\vergleichskette
{ Q
}
{ \subseteq }{ \R^n
}
{  }{
}
{    }{
}
{   }{
}
}
{}{}{}
ein achsenparalleler $n$-dimensionaler Quader
\zusatzklammer {mit Seiten aber ohne Kanten} {} {}
mit dem
\definitionsverweis {Rand}{}{}

\mathl{\partial Q}{} und $\omega$ eine auf $Q$ definierte
\definitionsverweis {stetig-differenzierbare}{}{}
$(n-1)$-\definitionsverweis {Differentialform}{}{.}}
\faktfolgerung {Dann ist

\mavergleichskettedisp
{\vergleichskette
{
\int_{  Q  }  d \omega
}
{ =} {
\int_{  \partial Q  }   \omega
}
{ } {
}
{ } {
}
{ } {
}
}
{}{}{.}}
\faktzusatz {}
\faktzusatz {}}

 }

__NOEDITSECTION__

\inputaufgabeklausurloesung
{0}

{

}
{  }

__NOEDITSECTION__

\inputaufgabeklausurloesung
{4 (1+3)}

{

Es sei $M$ die durch

\mavergleichskettedisp
{\vergleichskette
{ x^2+3y^2
}
{ =} { 2
}
{ } {
}
{ } {
}
{ } {
}
}
{}{}{}
gegebene Ellipse im $\R^2$,

\mavergleichskette
{\vergleichskette
{ P
}
{ = }{ \begin{pmatrix}  1  \\ { \frac{   1  }{  \sqrt{3}  } }    \end{pmatrix}
}
{  }{
}
{    }{
}
{   }{
}
}
{}{}{}
und

\mavergleichskette
{\vergleichskette
{ Q
}
{ = }{ \begin{pmatrix}  -1 \\ { \frac{   1  }{  \sqrt{3}  } }    \end{pmatrix}
}
{  }{
}
{    }{
}
{   }{
}
}
{}{}{}
Punkte auf $M$.
\aufzaehlungdrei{Zeige, dass

\mavergleichskettedisp
{\vergleichskette
{ v
}
{ =} { \begin{pmatrix}  -  \sqrt{3}  \\ 1    \end{pmatrix}
}
{ } {
}
{ } {
}
{ } {
}
}
{}{}{}
ein
\definitionsverweis {Tangentialvektor}{}{}
in $P$ an $M$  ist.
}{Bestimme das Bild von $v$ unter einem Paralleltransport auf $M$ von $P$ nach $Q$.
}{
}

}
{

\aufzaehlungzwei {Der Tangentialraum in einem Punkt $\begin{pmatrix}  x  \\ y   \end{pmatrix}$ steht senkrecht auf dem Gradienten
\mathl{\begin{pmatrix}  2x \\ 6y   \end{pmatrix}}{,} in $P$ also senkrecht auf

\mavergleichskettedisp
{\vergleichskette
{ \begin{pmatrix}  2 \\ { \frac{   6 }{  \sqrt{3}  } }    \end{pmatrix}
}
{ =} { 2 \begin{pmatrix}  1  \\ \sqrt{3}    \end{pmatrix}
}
{ } {
}
{ } {
}
{ } {
}
}
{}{}{.}
Dies erfüllt der Vektor $v$.
} {Der Paralleltransport längs eines Weges $\gamma$ auf $M$ von $P$ nach $Q$ ist eine Isometrie
\maabb {} {T_PM} {T_QM
} {.}
Der Tangentialraum an $Q$ steht senkrecht auf
\mathl{\begin{pmatrix}  - 1  \\ \sqrt{3}    \end{pmatrix}}{} und wird erzeugt von
\mathl{\begin{pmatrix}  \sqrt{3}  \\ 1    \end{pmatrix}}{.} Da sich die Norm nicht ändert, kommt als Ergebnis nur
\mathl{\pm \begin{pmatrix}  \sqrt{3}  \\ 1    \end{pmatrix}}{} in Frage. Da $v$ gegen den Uhrzeigersinn auf $M$ zeigt, muss dies auch für das Ergebnis gelten
\zusatzklammer {etwa wegen
[[Hyperfläche/Geodätische Kurve/Tangentiale Richtung/Paralleltransport/Aufgabe|Aufgabe 6.9 (Differentialgeometrie (Osnabrück 2023))]]} {} {,}
also ist
\mathl{- \begin{pmatrix}  \sqrt{3}  \\ 1    \end{pmatrix}}{} das Ergebnis.
}

 }

__NOEDITSECTION__

\inputaufgabeklausurloesung
{5}

{

Beweise den Satz über die Orientierungen auf einer differenzierbaren Hyperfläche.

}
{

Nach
[[Hyperfläche/Glatt/Normalenfeld/Fakt|Lemma 2.4 (Differentialgeometrie (Osnabrück 2023))]]
gibt es ein Einheitsnormalenfeld, das eine Orientierung repräsentiert, und die Negation davon liefert eine weitere Orientierung. Diese nennen wir $G$ bzw. $-G$. Es sei nun
\maabbdisp {N} { Y } { \R^n
} {}
ein beliebiges stetiges Einheitsnormalenfeld. Wir betrachten die Übereinstimmungsorte

\mavergleichskettedisp
{\vergleichskette
{ Y_1
}
{ =} { \left\{ P \in Y  \mid   N(P) = G(P)         \right\}
}
{ } {
}
{ } {
}
{ } {
}
}
{}{}{}
und

\mavergleichskettedisp
{\vergleichskette
{ Y_2
}
{ =} { \left\{ P \in Y  \mid   N(P) = - G(P)         \right\}
}
{ } {
}
{ } {
}
{ } {
}
}
{}{}{.}
Da es für jeden Punkt

\mavergleichskette
{\vergleichskette
{ P
}
{ \in }{ Y
}
{  }{
}
{    }{
}
{   }{
}
}
{}{}{}
nur die beiden möglichen Werte

\mavergleichskettedisp
{\vergleichskette
{ N(P)
}
{ =} { \pm G(P)
}
{ } {
}
{ } {
}
{ } {
}
}
{}{}{}
gibt, ist $Y$ die disjunkte Vereinigung von
\mathkor {} {Y_1} {und} {Y_2} {.}
Als Übereinstimmungsort von stetigen Funktionen sind
\mathkor {} {Y_1} {und} {Y_2} {}
\definitionsverweis {abgeschlossen}{}{}
\zusatzklammer {und damit auch
\definitionsverweis {offen}{}{}
in $Y$} {} {.}
Aufgrund des Zusammenhangs ist also

\mavergleichskette
{\vergleichskette
{ Y
}
{ = }{ Y_1
}
{  }{
}
{    }{
}
{   }{
}
}
{}{}{}
oder

\mavergleichskette
{\vergleichskette
{ Y
}
{ = }{ Y_2
}
{  }{
}
{    }{
}
{   }{
}
}
{}{}{.}

 }

__NOEDITSECTION__

\inputaufgabeklausurloesung
{4}

{

Es sei

\mavergleichskette
{\vergleichskette
{ Y
}
{ \subseteq }{ W
}
{ \subseteq }{ \R^3
}
{    }{
}
{   }{
}
}
{}{}{,}
$W$
\definitionsverweis {offen}{}{,}
eine
\definitionsverweis {differenzierbare Fläche}{}{}
versehen mit einem
\definitionsverweis {Einheitsnormalenfeld}{}{}
$N$. Es sei
\maabb {\gamma} { I } { Y
} {}
eine stetig differenzierbare Kurve und sei
\maabbdisp {F} { I } { \R^3
} {}
ein differenzierbares
\definitionsverweis {tangentiales Vektorfeld}{}{}
längs $\gamma$, das
\definitionsverweis {parallel}{}{}
sei. Zeige, dass dann auch das Vektorfeld
\maabbeledisp {} { I } { \R^3
} { t } { N(\gamma(t)) \times F(t)
} {,}
parallel ist, wobei $\times$ das
\definitionsverweis {Kreuzprodukt}{}{}
im $\R^3$ bezeichnet.

}
{

Es ist

\mavergleichskettealign
{\vergleichskettealign
{ (N (\gamma(t)) \times F(t))'
}
{ =} { (N (\gamma(t)))' \times F(t) +  N (\gamma(t)) \times F'(t)
}
{ =} { \left(  D N   \right)_{ \gamma(t)} \left(   \gamma'(t)   \right)  \times F(t) +  N (\gamma(t)) \times F'(t)
}
{ } {
}
{ } {
}
}
{}
{}{.}
Es ist zu zeigen, dass dieser Vektor stets senkrecht auf dem Tangentialraum $T_{\gamma(t)}Y$ steht. Da $F'(t)$ senkrecht auf dem
Tangentialraum steht, ist er linear abhängig zum Normalenvektor $N (\gamma(t))$ und daher ist der rechte Summand gleich $0$ nach
[[K^3/Kreuzprodukt/Eigenschaften/Fakt|Lemma 33.3 (Lineare Algebra (Osnabrück 2024-2025))  (3)]].
Im linken Summanden ist $F(t)$ tangential nach Voraussetzung und $\left(  D N   \right)_{ \gamma(t)} \left(   \gamma'(t)   \right)$ ist tangential nach
[[Hyperfläche/Weingartenabbildung/Tangentialraum/Fakt|Lemma 4.2 (Differentialgeometrie (Osnabrück 2023))]].
Daher ist nach
[[K^3/Kreuzprodukt/Eigenschaften/Fakt|Lemma 33.3 (Lineare Algebra (Osnabrück 2024-2025))  (6)]]
der linke Summand senkrecht auf dem Tangentialraum und somit ist insgesamt das Vektorfeld parallel.

 }

__NOEDITSECTION__

\inputaufgabeklausurloesung
{9}

{

Es sei

\mavergleichskette
{\vergleichskette
{ M
}
{ \neq }{ \emptyset
}
{  }{
}
{    }{
}
{   }{
}
}
{}{}{}
eine
\definitionsverweis {differenzierbare Mannigfaltigkeit}{}{}
der Dimension $n$. Zeige, dass es eine Kette von
\definitionsverweis {abgeschlossenen Untermannigfaltigkeiten}{}{}

\mavergleichskettedisp
{\vergleichskette
{ M_0
}
{ \subseteq} { M_1
}
{ \subseteq} { M_2
}
{ \subseteq  \ldots \subseteq} { M_{n-1}
}
{ \subseteq} { M_n
}
}
{
\vergleichskettefortsetzung
{ =} { M
}
{ } {}
{ } {}
{ } {}
}{}{}
derart gibt, dass die abgeschlossene Untermannigfaltigkeit $M_i$ die Dimension $i$ besitzt.

}
{

Es sei $P \in M$ ein Punkt und
\mathl{P \in U}{} eine offene Kartenumgebung zusammen mit einer Karte
\maabbdisp {\alpha} { U } {V \subseteq \R^n
} {,}
wobei $V=B(0,3)$ die offene Kugel mit Mittelpunkt $0$ und Radius $3$ sei und wobei
\mathl{\alpha(P)=0}{} gelte. Für
\mathl{i=1 , \ldots , n-1}{}  betrachten wir

\mathdisp {B_i = \left\{  x \in V  \mid  (x_1-1)^2+x_2^2 + \cdots + x_n^2 = 1 , \,  x_{i+2}   = 0,\, x_{i+3} = 0 , \ldots , x_n = 0       \right\}} { . }

Es ist also $B_{n-1}$ die Kugeloberfläche mit Mittelpunkt
\mathl{(1,0 , \ldots , 0)}{} und Radius $1$, $B_{n-2}$ ist darin der durch
\mathl{x_n=0}{} definierte \anfuehrung{Äquator}{} usw. Man erhält $B_i$ aus
\mathl{B_{i+1}}{,} indem man zusätzlich noch
\mathl{x_{i+2}=0}{} setzt. Daher liegt eine absteigende Kette von abgeschlossenen Teilmengen

\mathdisp {B_{n-1} \supseteq B_{n-2} \supseteq  \ldots \supseteq B_1 \supseteq B_0} {  }
 vor
\zusatzklammer {$B_0$ besteht aus den beiden Punkten
\mathl{(\pm1,0 , \ldots , 0 )}{}} {} {.}
Wir fassen $B_i$ als die Faser über dem Nullpunkt der Abbildung
\maabbeledisp {\varphi_i} { V } {\R^{n-i}
} {(x_1 , \ldots , x_n) } {  ((x_1-1)^2+x_2^2 + \cdots + x_n^2 - 1 , x_{i+2} , \ldots , x_n)
} {,}
auf. Die Jacobimatrix dieser Abbildung ist

\mathdisp {\begin{pmatrix}  2x_1-2 &  2x_2  &  \ldots  &  2x_{i+1} &  2x_{i+2} &  \ldots &  2x_{n-1} &  2x_n
 \\  0 &  0 &  \ldots &  0 &  1 &  0 &  \ldots &  0
 \\ \vdots  & \vdots & \vdots &  \ddots &  \ddots &  \ddots  &  \ddots  & \vdots \\  0 &  0 &  \ldots &  \ldots &  \ldots &  0 &  1 &  0 \\  0 &  0 &  \ldots &  \ldots &  \ldots  &  \ldots &  0 &  1 \end{pmatrix}} { . }

Der Rang dieser Matrix ist nur bei
\mathl{x_1=1, \, x_2 = \cdots =x_{i+1}=0}{} kleiner als $n-i$, ein solcher Punkt liegt also nicht auf $B_i$. Das bedeutet, dass die Abbildung $\varphi_i$ in der Faser $B_i$ regulär ist, sodass $B_i$ aufgrund des Satzes über implizite Abbildungen eine abgeschlossene Untermannigfaltigkeit von $V$ der Dimension
\mathl{i}{} ist.

Wir setzen nun
\mathl{M_i= \alpha^{-1}(B_i)}{} für
\mathl{i=1 , \ldots , n-1}{} und
\mathl{M_n= M}{.} Da die $B_i$ kompakt sind, sind die $M_i$ auch abgeschlossene Teilmengen in $M$. Da die Bedingung für eine abgeschlossene Mannigfaltigkeit eine lokale Eigenschaft ist, handelt es sich um abgeschlossene Untermannigfaltigkeiten von $M$.

 }

__NOEDITSECTION__

\inputaufgabeklausurloesung
{0}

{

}
{  }

__NOEDITSECTION__

\inputaufgabeklausurloesung
{8}

{

Es sei
\maabbdisp {\psi} {\R^n } {\R
} {}
eine
\definitionsverweis {differenzierbare Funktion}{}{}
und

\mavergleichskettedisphandlinks
{\vergleichskettedisphandlinks
{M
}
{ =} { \left\{  \left(  x_1   , \,   , \ldots ,  , \,   x_n , \,   \psi(x_1 , \ldots , x_n)                \right)   \mid  (x_1 , \ldots , x_n) \in \R^n         \right\}
}
{ \subseteq} { \R^n \times \R
}
{ } {
}
{ } {
}
}
{}{}{}
der
\definitionsverweis {Graph}{}{}
von $\psi$. Zeige, dass $M$ eine
\definitionsverweis {orientierte}{}{} \definitionsverweis {riemannsche Mannigfaltigkeit}{}{}
ist und dass für die
\definitionsverweis {kanonische Volumenform}{}{}
$\omega$ auf $M$ die Formel

\mavergleichskettedisp
{\vergleichskette
{ (\operatorname{Id} \times \psi)^*(\omega)
}
{ =} { \left(  1+ \sum_{i = 1}^n \left(  { \frac{   \partial   \psi  }{  \partial   x_i   } }  \right)^2  \right)^{1/2} dx_1 \wedge \ldots \wedge dx_n
}
{ } {
}
{ } {
}
{ } {
}
}
{}{}{}
gilt.

}
{

Es sei

\mavergleichskette
{\vergleichskette
{ W
}
{ = }{ \R^n
}
{  }{
}
{    }{
}
{   }{
}
}
{}{}{.}
Die Abbildung
\maabbeledisp {\varphi = \operatorname{id} \times \psi} { W } { W \times \R
} { x } {(x, \psi(x))
} {,}
ist ein
\definitionsverweis {Diffeomorphismus}{}{}
zwischen $W$ und dem Graphen $M$. Der Graph ist eine
\definitionsverweis {abgeschlossene Untermannigfaltigkeit}{}{}
von
\mathl{W \times \R}{} und trägt daher die
\definitionsverweis {induzierte riemannsche Struktur}{}{}
und
\zusatzklammer {da sich die Orientierung von $W$ auf $M$ überträgt} {} {}
eine
\definitionsverweis {kanonische Volumenform}{}{}
$\omega$. Auf diese Situation kann man
[[Riemannsche Untermannigfaltigkeit des R^n/Einbettung/Volumenform/Fakt|Satz 16.8 (Differentialgeometrie (Osnabrück 2023))]]
anwenden. Die
\definitionsverweis {partiellen Ableitungen}{}{}
von $\varphi$ nach der $i$-ten Variablen sind

\mavergleichskettedisp
{\vergleichskette
{ { \frac{   \partial  \varphi  }{  \partial   x_i   } }
}
{ =} { \begin{pmatrix}  0  \\\vdots\\  0\\ 1\\  0\\ \vdots\\  0\\  { \frac{   \partial   \psi  }{  \partial   x_i   } }   \end{pmatrix}
}
{ } {
}
{ } {
}
{ } {
}
}
{}{}{.}
Es sei

\mavergleichskette
{\vergleichskette
{Q
}
{ \in }{ W
}
{  }{
}
{    }{
}
{   }{
}
}
{}{}{}
ein Punkt, den wir in die Funktionen im Folgenden einsetzen, sodass wir überall mit reellen Zahlen rechnen. Die Skalarprodukte, die die Einträge
\mathl{b_{ij}}{} der Matrix $B$ bilden
\zusatzklammer {von deren Determinante wir die Wurzel berechnen müssen} {} {,}
sind gleich

\mavergleichskettedisp
{\vergleichskette
{ b_{ij}
}
{ =} { \left\langle  { \frac{   \partial  \varphi  }{  \partial   x_i   } }(Q)  ,  { \frac{   \partial  \varphi  }{  \partial   x_j   } }(Q)   \right\rangle
}
{ } {}
{ } {}
{ } {}
}
{}{}{}
Wir schreiben

\mavergleichskette
{\vergleichskette
{B
}
{ = }{ E_n+A
}
{  }{
}
{    }{
}
{   }{
}
}
{}{}{}
mit

\mavergleichskette
{\vergleichskette
{ A
}
{ = }{ \left(  a_{ij}  \right)_{1 \leq i,j \leq n}
}
{  }{
}
{    }{
}
{   }{
}
}
{}{}{.}
Mit

\mavergleichskette
{\vergleichskette
{ c_i
}
{ = }{ { \frac{   \partial   \psi  }{  \partial   x_i   } } ( Q)
}
{  }{
}
{    }{
}
{   }{
}
}
{}{}{}
können wir

\mavergleichskette
{\vergleichskette
{ a_{ij}
}
{ = }{ c_ic_j
}
{  }{
}
{    }{
}
{   }{
}
}
{}{}{}
und insgesamt die Matrix $A$ als

\mavergleichskettedisp
{\vergleichskette
{A
}
{ =} { \begin{pmatrix} c_1  \\\vdots\\ c_n   \end{pmatrix} \left( c_1   , \,   \ldots  , \,  c_n                 \right)
}
{ } {
}
{ } {
}
{ } {
}
}
{}{}{}
schreiben. Daher beschreibt $A$ eine lineare Abbildung von $\R^n$ nach $\R^n$, die durch $\R$
\definitionsverweis {faktorisiert}{}{,}
und besitzt damit einen
\definitionsverweis {Kern}{}{,}
der zumindest $(n-1)$-dimensional ist. Nennen wir ihn $K$. Wenn er die Dimension $n$ besitzt, so ist

\mavergleichskette
{\vergleichskette
{A
}
{ = }{ 0
}
{  }{
}
{    }{
}
{   }{
}
}
{}{}{}
und $B$ ist die Identität, und die Aussage ist richtig. Es sei also

\mavergleichskette
{\vergleichskette
{A
}
{ \neq }{ 0
}
{  }{
}
{    }{
}
{   }{
}
}
{}{}{.}
Dann ist

\mavergleichskette
{\vergleichskette
{v
}
{ = }{ \begin{pmatrix} c_1  \\\vdots\\ c_n   \end{pmatrix}
}
{  }{
}
{    }{
}
{   }{
}
}
{}{}{}
ein
\definitionsverweis {Eigenvektor}{}{}
von $A$ zum
\definitionsverweis {Eigenwert}{}{}

\mavergleichskette
{\vergleichskette
{ c_1^2 + \cdots + c_n^2
}
{ \neq }{ 0
}
{  }{
}
{    }{
}
{   }{
}
}
{}{}{.}
Dieser Vektor ist ein Eigenvektor von $B$ zum Eigenwert
\mathl{1+c_1^2 + \cdots + c_n^2}{} und $K$ bildet den
\mathl{(n-1)}{-}dimensionalen
\definitionsverweis {Eigenraum}{}{}
für $B$ zum Eigenwert $1$. Insgesamt ist $B$
\definitionsverweis {diagonalisierbar}{}{}
und ihre
\definitionsverweis {Determinante}{}{}
ist das Produkt der Eigenwerte, also gleich
\mathl{1+c_1^2 + \cdots + c_n^2}{.}

 }

__NOEDITSECTION__

\inputaufgabeklausurloesung
{6}

{

Berechne die Oberfläche der Einheitskugel.

}
{  }

__NOEDITSECTION__

\inputaufgabeklausurloesung
{0}

{

}
{  }

__NOEDITSECTION__

\inputaufgabeklausurloesung
{3}

{

Berechne die
\definitionsverweis {äußere Ableitung}{}{}
$d \omega$ der
$1$-\definitionsverweis {Differentialform}{}{}

\mavergleichskettedisp
{\vergleichskette
{ \omega
}
{ =} { { \frac{   x^2 }{  y } } dx- { \frac{   x  }{  y^2 } } dy
}
{ } {
}
{ } {
}
{ } {
}
}
{}{}{}
auf

\mavergleichskette
{\vergleichskette
{ U
}
{ = }{ \left\{ (x,y) \in \R^2  \mid   y \neq 0         \right\}
}
{  }{
}
{    }{
}
{   }{
}
}
{}{}{.}

}
{

Es ist

\mavergleichskettealign
{\vergleichskettealign
{d \omega
}
{ =} { d { \frac{   x^2 }{  y } } \wedge dx - d { \frac{   x  }{  y^2 } } \wedge dy
}
{ =} { - { \frac{   x^2 }{  y^2 } } dy  \wedge dx  -   { \frac{   1 }{  y^2 } }  dx \wedge dy
}
{ =} { { \frac{   x^2 -1 }{  y^2 } }  dx \wedge dy
}
{ } {
}
}
{}
{}{.}

 }

__NOEDITSECTION__

\inputaufgabeklausurloesung
{3}

{

Man gebe ein Beispiel für eine
\definitionsverweis {zusammenhängende}{}{}
\definitionsverweis {Mannigfaltigkeit mit Rand}{}{,}
deren Rand aus
\zusatzklammer {einer disjunkten Vereinigung von} {} {}
unendlich vielen Kreisen $S^1$ besteht.

}
{

Wir berachten

\mavergleichskettedisp
{\vergleichskette
{ M
}
{ =} { \R^2 \setminus \biguplus_{n \in \Z} U \left(  (n,0), { \frac{   1 }{  3 } }   \right)
}
{ } {
}
{ } {
}
{ } {
}
}
{}{}{.}
Es werden also aus der reellen Ebene längs der $x$-Achse unendlich viele disjunkte offene Bälle herausgenommen. Dadurch entsteht zu jedem Ball ein Teilrand von $M$, der diffeomorph zur Kreissphäre $S^1$ ist. $M$ ist offenbar wegzusammenhängend.

 }

__NOEDITSECTION__

\inputaufgabeklausurloesung
{0}

{

}
{  }

__NOEDITSECTION__

\inputaufgabeklausurloesung
{0}

{

}
{  }

__NOEDITSECTION__

\inputaufgabeklausurloesung
{7}

{

Es sei $M$ eine
\definitionsverweis {riemannsche Mannigfaltigkeit}{}{}
und sei ein
\definitionsverweis {linearer Zusammenhang}{}{}
auf dem
\definitionsverweis {Tangentialbündel}{}{}
$TM$ gegeben, der
\definitionsverweis {metrisch}{}{}
und
\definitionsverweis {torsionsfrei}{}{}
sei. Zeige, dass für Vektorfelder $X,Y,Z$ die Formel

\mavergleichskettedisp
{\vergleichskette
{ \left\langle \nabla_X Y , Z  \right\rangle
}
{ =} { { \frac{   1  }{  2 } }  \left(  D_X \left\langle  Y  ,  Z  \right\rangle +  D_Y \left\langle  Z  ,  X  \right\rangle -  D_Z \left\langle  X  ,  Y  \right\rangle  - \left\langle  Y  , [X,Z]   \right\rangle - \left\langle  Z  , [Y,X]   \right\rangle + \left\langle  X  , [Z,Y]   \right\rangle   \right)
}
{ } {
}
{ } {
}
{ } {
}
}
{}{}{}
gilt.

}
{

Wegen der Torsionsfreiheit gilt

\mavergleichskettedisp
{\vergleichskette
{ \nabla_VW- \nabla_WV
}
{ =} { [V,W]
}
{ } {
}
{ } {
}
{ } {
}
}
{}{}{,}

\mavergleichskettedisp
{\vergleichskette
{ \nabla_VZ- \nabla_ZV
}
{ =} { [V,Z]
}
{ } {
}
{ } {
}
{ } {
}
}
{}{}{}
und

\mavergleichskettedisp
{\vergleichskette
{ \nabla_WZ- \nabla_ZW
}
{ =} { [W,Z]
}
{ } {
}
{ } {
}
{ } {
}
}
{}{}{,}
wobei wir die erste Identität auch als

\mavergleichskettedisp
{\vergleichskette
{ \nabla_VW + \nabla_WV
}
{ =} { 2 \nabla_VW -[V,W]
}
{ } {
}
{ } {
}
{ } {
}
}
{}{}{}
auffassen. Wegen metrisch gelten die Identitäten

\mavergleichskettedisp
{\vergleichskette
{ D_V \left\langle  W  ,  Z  \right\rangle
}
{ =} { \left\langle \nabla_V W , Z  \right\rangle  + \left\langle  W  ,  \nabla_V Z  \right\rangle
}
{ } {
}
{ } {
}
{ } {
}
}
{}{}{,}

\mavergleichskettedisp
{\vergleichskette
{ D_W \left\langle  Z  ,  V  \right\rangle
}
{ =} { \left\langle \nabla_WZ , V  \right\rangle  + \left\langle  Z  ,  \nabla_WV  \right\rangle
}
{ } {
}
{ } {
}
{ } {
}
}
{}{}{}
und

\mavergleichskettedisp
{\vergleichskette
{ D_Z \left\langle  V  ,  W  \right\rangle
}
{ =} { \left\langle \nabla_Z V , W  \right\rangle  + \left\langle  V  ,  \nabla_Z W   \right\rangle
}
{ } {
}
{ } {
}
{ } {
}
}
{}{}{.}
Wir addieren die Gleichungen zusammen und erhalten

\mavergleichskettedisphandlinks
{\vergleichskettedisphandlinks
{ D_V \left\langle  W  ,  Z  \right\rangle+ D_W \left\langle  Z  ,  V  \right\rangle  - D_Z \left\langle  V  ,  W  \right\rangle
}
{ =} { \left\langle \nabla_V W+\nabla_W V  ,  Z  \right\rangle + \left\langle \nabla_W Z-\nabla_Z W   ,  V  \right\rangle+\left\langle  \nabla_V Z-\nabla_Z V   ,  W  \right\rangle
}
{ } {
}
{ } {
}
{ } {
}
}
{}{}{.}
In die rechte Seite setzen wir die oben erzielten Ausdrücke ein und erhalten

\mavergleichskettealignhandlinks
{\vergleichskettealignhandlinks
{ D_V \left\langle  W  ,  Z  \right\rangle+ D_W \left\langle  Z  ,  V  \right\rangle  - D_Z \left\langle  V  ,  W  \right\rangle
}
{ =} { \left\langle  2 \nabla_VW -[V,W]  ,  Z  \right\rangle + \left\langle  [W,Z]  ,  V  \right\rangle + \left\langle [V,Z]   ,  W  \right\rangle
}
{ =} { 2\left\langle  \nabla_VW   ,  Z  \right\rangle - \left\langle  [V,W]  ,  Z  \right\rangle + \left\langle  [W,Z]  ,  V  \right\rangle + \left\langle  [V,Z]   ,  W  \right\rangle
}
{ } {
}
{ } {
}
}
{}
{}{.}
Eine Umstellung ergibt die Formel.

Aufgrund der Gleichung ist
\mathl{\left\langle  \nabla_V W , Z  \right\rangle}{} für beliebige Vektorfelder durch die rechte Seite festgelegt, in der der Zusammenhang gar nicht vorkommt. Da dies für jedes Vektorfeld $Z$ gilt, ist dadurch auch die vertikale Ableitung $\nabla_V W$ und damit der lineare Zusammenhang eindeutig festgelegt.

 }

 [[Kategorie:Latexseite]]
